\documentclass[11pt,letterpaper]{article}

% --- Packages ---
\usepackage[margin=1in, headheight=24pt]{geometry}
\usepackage{graphicx}
\usepackage{hyperref}
\usepackage{enumitem}
\usepackage{booktabs}
\usepackage{xcolor}
\usepackage{tcolorbox}
\usepackage{fancyhdr}
\usepackage{titlesec}
\usepackage{parskip}
\usepackage[T1]{fontenc}
\usepackage{microtype}
% \usepackage[utf8]{inputenc}  % UTF-8 is default in modern LaTeX
\usepackage{amssymb}
\usepackage{amsmath}
\usepackage{tabularx}

% --- Colors ---
\definecolor{minesblue}{RGB}{33,69,126}
\definecolor{minesgold}{RGB}{204,174,0}
\definecolor{tipgreen}{RGB}{34,120,58}
\definecolor{warnred}{RGB}{180,30,30}
\definecolor{bglight}{RGB}{245,245,245}

% --- Hyperref Setup ---
\hypersetup{
    colorlinks=true,
    linkcolor=minesblue,
    urlcolor=minesblue,
    citecolor=minesblue,
    pdftitle={MEGN 300 Student Guide: FAQ, Tips, and Common Pitfalls},
    pdfauthor={Adam Duran},
    pdfsubject={MEGN 300 Instrumentation and Automation},
}

% --- Header/Footer ---
\pagestyle{fancy}
\fancyhf{}
\fancyhead[L]{\textcolor{minesblue}{\textbf{MEGN 300 -- Instrumentation and Automation}}}
\fancyhead[R]{\textcolor{minesblue}{Spring 2026}}
\fancyfoot[C]{\thepage}
\renewcommand{\headrulewidth}{1pt}
\renewcommand{\headrule}{\hbox to\headwidth{\color{minesblue}\leaders\hrule height \headrulewidth\hfill}}

% --- Custom tcolorbox environments ---
\newtcolorbox{faqbox}[1][]{colback=bglight,colframe=minesblue,fonttitle=\bfseries,title={#1},arc=2pt,boxrule=0.8pt,left=6pt,right=6pt,top=4pt,bottom=4pt}
\newtcolorbox{tipbox}{colback=tipgreen!5,colframe=tipgreen,fonttitle=\bfseries\color{tipgreen},title={Tip},arc=2pt,boxrule=0.8pt,left=6pt,right=6pt,top=4pt,bottom=4pt}
\newtcolorbox{warnbox}{colback=warnred!5,colframe=warnred,fonttitle=\bfseries\color{warnred},title={Common Pitfall},arc=2pt,boxrule=0.8pt,left=6pt,right=6pt,top=4pt,bottom=4pt}
\newtcolorbox{protip}{colback=minesgold!8,colframe=minesgold!80!black,fonttitle=\bfseries\color{minesgold!80!black},title={Pro Tip},arc=2pt,boxrule=0.8pt,left=6pt,right=6pt,top=4pt,bottom=4pt}
\newtcolorbox{priorities}{colback=minesblue!5,colframe=minesblue!60,fonttitle=\bfseries\color{minesblue},title={\small$\blacktriangleright$ Top Priorities This Module},arc=2pt,boxrule=0.6pt,left=6pt,right=6pt,top=3pt,bottom=3pt,before upper={\small}}

% --- Section formatting ---
\titleformat{\section}
    {\Large\bfseries\color{minesblue}}
    {\thesection}{1em}{}[\vspace{-4pt}\textcolor{minesgold}{\rule{\textwidth}{1.5pt}}]
\titleformat{\subsection}
    {\large\bfseries\color{minesblue!80}}
    {\thesubsection}{1em}{}

% --- Document ---
\begin{document}

% --- Title ---
\begin{center}
    {\LARGE\bfseries\textcolor{minesblue}{MEGN 300 Instrumentation \& Automation}}\\[6pt]
    {\Large\textcolor{minesblue!70}{Student Guide: FAQ, Tips, and Common Pitfalls}}\\[8pt]
    {\large Spring 2026}\\[4pt]
    \textcolor{minesgold}{\rule{0.6\textwidth}{2pt}}
\end{center}

\vspace{0.5cm}

\noindent This guide accompanies the \textit{MEGN~300 Master Reference Document} and the official course Canvas page. It is \textbf{not a replacement} for the syllabus, lab handouts, or Canvas; always refer to those for exact deadlines and submission requirements. The purpose of this guide is to help you anticipate common questions, avoid frequent mistakes, and build strong instrumentation habits. Read the relevant section \textit{before} each lab module.

\begin{tcolorbox}[colback=minesgold!10, colframe=minesgold!70!black, boxrule=0.6pt, arc=2pt, left=6pt, right=6pt, top=4pt, bottom=4pt]
\textbf{How to use this guide:} Before each lab or design challenge, open this guide to the relevant module section. Read the \textit{Common Pitfall} and \textit{Pro Tip} boxes first; those address the mistakes that cost teams the most time and points.
\end{tcolorbox}

\begin{warnbox}
\textbf{Read the AI Use Policy in the MRD before your first assignment.} The Master Reference Document contains a ``Modern Tools and AI Use Policy'' chapter with a Green/Yellow/Red traffic-light framework. AI tools are permitted for learning (Green) and with disclosure for scaffolding (Yellow), but submitting AI-generated work as your own is an academic integrity violation (Red). Apply the Mastery Test: can you do it without the AI?
\end{warnbox}

\tableofcontents
\newpage

%=============================================================================
\section{Course Philosophy}
%=============================================================================

\begin{faqbox}[Q: What is MEGN 300 really about?]
MEGN~300 teaches you to \textbf{measure physical quantities accurately and control engineering systems automatically}. It bridges theory (physics, math, statistics) and practice (circuits, sensors, LabVIEW, NI~DAQ hardware). Every concept leads to a hands-on skill: if you understand error analysis, you can design a measurement system with quantified uncertainty. If you understand PID control, you can regulate temperature, speed, or pressure in a real system.
\end{faqbox}

\begin{faqbox}[Q: What resources should I be using?]
Three documents work together:
\begin{enumerate}[leftmargin=*, itemsep=1pt, topsep=0pt]
    \item \textbf{This Student Guide}: quick-hit tips, common pitfalls, and FAQ for each module.
    \item \textbf{The Master Reference Document}: deep technical content organized in five parts (Measurement Fundamentals, Signals \& Signal Processing, Electronics including debugging methodology, Control Systems, and Engineering Science Reference). Every chapter has worked ThermoRig examples, practice questions with solutions, and Requirements Mapping boxes. Consult it when you need the equations, theory, or design procedures. Start with the ``What Should I Be Reading Right Now?'' triage guide on page~2 to find exactly the section you need.
    \item \textbf{Canvas}: assignment details, deadlines, rubrics, lab handouts, and LabVIEW starter VIs.
\end{enumerate}
\end{faqbox}

\begin{protip}
\textbf{The Master Reference Document has a LabVIEW appendix with a Signal Walk Debugging Protocol.} Before every lab, check the appendix for the DAQmx programming pattern, essential VIs, and the ``Common Mistakes'' checklist. Most LabVIEW grading deductions come from the six mistakes listed there.
\end{protip}

%=============================================================================
\section{Module 1: Measurement Fundamentals and Static Calibration}
%=============================================================================

\begin{priorities}
\begin{enumerate}[leftmargin=*, itemsep=1pt, topsep=0pt]
    \item Understand the five-block measurement system and where errors enter.
    \item Know the difference between accuracy, precision, systematic, and random error.
    \item Perform a static calibration: collect data, fit a line, and report uncertainty.
\end{enumerate}
\end{priorities}

\begin{tcolorbox}[colback=minesblue!3, colframe=minesblue!40, boxrule=0.5pt, arc=2pt, left=6pt, right=6pt, top=3pt, bottom=3pt, fonttitle=\small\bfseries\color{minesblue}, title={Before You Start: Module 1}]
\small
\begin{enumerate}[nosep,leftmargin=*]
\item Read MRD Chapter~1 (Fundamentals of Measurement Theory) and Chapter~2 (Static Measurements).
\item Review the ``First Week Survival Guide'' in the MRD front matter.
\item Watch the LabVIEW DAQmx tutorial video on Canvas.
\item Familiarize yourself with the lab bench: locate the NI~USB-6001, bench supply, DMM, and component drawers.
\end{enumerate}
\end{tcolorbox}

\begin{faqbox}[Q: What is the difference between accuracy and precision?]
\textbf{Accuracy} is closeness to the true value (affected by systematic/bias errors). \textbf{Precision} is closeness of repeated measurements to each other (affected by random errors). A sensor can be precise but inaccurate (all readings cluster together but off-center). Calibration fixes accuracy. Averaging improves precision. You need both. The Master Reference Document, Chapter~1 (Fundamentals of Measurement Theory), has the dartboard analogy and worked examples.
\end{faqbox}

\begin{warnbox}
\textbf{The \#1 calibration mistake}: not waiting for the sensor to reach steady state at each calibration point. If you are calibrating a temperature sensor and change the bath temperature, wait at least $5\tau$ (five time constants) before recording. A thermocouple with $\tau = 1$\,s needs 5\,s; an RTD probe with $\tau = 10$\,s needs 50\,s. Recording during the transient produces systematically wrong calibration data.
\end{warnbox}

\begin{tipbox}
When writing your calibration report, always report: the calibration equation (with coefficients), the range over which it is valid, the maximum residual, the number of calibration points, and the uncertainty of your reference standard. A calibration without documented uncertainty is just a guess.
\end{tipbox}

\begin{protip}
\textbf{Learn the Wheatstone bridge now; it appears everywhere.} Strain gauges, RTDs, pressure sensors, and load cells all use bridge circuits. The Master Reference Document, Chapter~2 (Static Measurements), walks through quarter-, half-, and full-bridge configurations with the RTD ThermoRig example and self-heating calculations.
\end{protip}

%=============================================================================
\section{Module 2: Dynamic Measurements and Time Constants}
%=============================================================================

\begin{priorities}
\begin{enumerate}[leftmargin=*, itemsep=1pt, topsep=0pt]
    \item Measure $\tau$ from a step response (first-order sensor).
    \item Understand the bandwidth-rise time relationship: $f_{-3\text{dB}} \times t_r \approx 0.35$ (for a first-order system).
    \item Know when a sensor's dynamic response will distort your measurement.
\end{enumerate}
\end{priorities}

\begin{tcolorbox}[colback=minesblue!3, colframe=minesblue!40, boxrule=0.5pt, arc=2pt, left=6pt, right=6pt, top=3pt, bottom=3pt, fonttitle=\small\bfseries\color{minesblue}, title={Before You Start: Module 2}]
\small
\begin{enumerate}[nosep,leftmargin=*]
\item Read MRD Chapter~3 (Dynamic Measurements). Focus on the step response figure and the Rule of Five.
\item Review the $5\tau$ derivation: $e^{-5} = 0.0067$, so $y(5\tau)/y_\infty = 99.3\%$.
\item Prepare to measure a thermocouple step response in lab: boiling water, ice water, and a stopwatch.
\end{enumerate}
\end{tcolorbox}

\begin{faqbox}[Q: How do I measure the time constant of a sensor?]
Apply a step change (e.g., plunge a thermocouple from room temperature into boiling water). Record the output vs.\ time. The time constant $\tau$ is the time to reach 63.2\% of the final value. Or, plot $\ln(1 - y/y_\infty)$ vs.\ $t$; the slope is $-1/\tau$. See the Master Reference Document, Chapter~3 (Dynamic Measurements), for the full procedure and the ThermoRig worked example.
\end{faqbox}

\begin{warnbox}
\textbf{Forgetting that sensors have bandwidth limits.} A thermocouple with $\tau = 0.2$\,s cannot faithfully measure a temperature oscillating at 1\,Hz; it will attenuate the signal to 62\% and add 52\textdegree{} of phase lag. If your measurement is dynamic, always check that the sensor bandwidth exceeds $5\times$ your signal frequency. The Master Reference Document has the attenuation formulas.
\end{warnbox}

\begin{tipbox}
\textbf{The Rule of Five}: wait $5\tau$ for settling. Use $f_{\text{sensor}} \geq 5 f_{\text{signal}}$ for flat response. Both give $<$1\% error.
\end{tipbox}

%=============================================================================
\section{Module 3: Error Analysis and Uncertainty}
%=============================================================================

\begin{priorities}
\begin{enumerate}[leftmargin=*, itemsep=1pt, topsep=0pt]
    \item Calculate mean, standard deviation, standard error, and 95\% confidence intervals.
    \item Propagate uncertainty through a multi-variable equation using RSS.
    \item Report measurements correctly: $\bar{x} \pm U$ (95\% CI, $n$ = ...).
\end{enumerate}
\end{priorities}

\begin{tcolorbox}[colback=minesblue!3, colframe=minesblue!40, boxrule=0.5pt, arc=2pt, left=6pt, right=6pt, top=3pt, bottom=3pt, fonttitle=\small\bfseries\color{minesblue}, title={Before You Start: Module 3}]
\small
\begin{enumerate}[nosep,leftmargin=*]
\item Read MRD Chapter~4 (Error Analysis and Uncertainty).
\item Review the RSS propagation formula and work through the ThermoRig heat flux example.
\item Bring a calculator to lab; you will compute uncertainty in real time.
\end{enumerate}
\end{tcolorbox}

\begin{faqbox}[Q: When do I use $n$ vs.\ $n-1$ in the standard deviation?]
Use $n-1$ (Bessel's correction) for the \textbf{sample} standard deviation when your measurements are a sample of a larger population (this is almost always the case in MEGN~300). Use $n$ only when you have measured the \emph{entire} population.
\end{faqbox}

\begin{warnbox}
\textbf{You cannot average away systematic error.} If your instrument has a 2\,\textdegree{}C calibration offset, averaging 1000 readings gives you a very precise number that is still 2\,\textdegree{}C wrong. Calibrate first, then average. The Master Reference Document, Chapter~4 (Error Analysis), has a warning box dedicated to this.
\end{warnbox}

\begin{faqbox}[Q: What is the RSS uncertainty propagation formula?]
For a result $f(x_1, x_2, \ldots, x_n)$ with independent input uncertainties $u_{x_i}$:
\[
u_f = \sqrt{\left(\frac{\partial f}{\partial x_1}\right)^2 u_{x_1}^2 + \left(\frac{\partial f}{\partial x_2}\right)^2 u_{x_2}^2 + \cdots + \left(\frac{\partial f}{\partial x_n}\right)^2 u_{x_n}^2}
\]
Each $\partial f / \partial x_i$ is the sensitivity coefficient. The largest term dominates the total.
\end{faqbox}

\begin{protip}
\textbf{After propagating uncertainty, identify the dominant contributor.} In the RSS formula, the term with the largest partial derivative times its uncertainty dominates the total. Improving that single measurement improves the overall result more than improving all others combined. The ThermoRig heat flux example in Chapter~4 (Error Analysis) shows how thickness ($L$) dominates even though it seems like the simplest measurement.
\end{protip}

%=============================================================================
\section{Module 4: Analog vs.\ Digital, Sampling, and Aliasing}
%=============================================================================

\begin{priorities}
\begin{enumerate}[leftmargin=*, itemsep=1pt, topsep=0pt]
    \item Understand the Nyquist criterion: $f_s > 2 f_{\max}$.
    \item Know that aliasing cannot be undone after sampling.
    \item Calculate ADC resolution (LSB) for a given bit depth and range.
\end{enumerate}
\end{priorities}

\begin{tcolorbox}[colback=minesblue!3, colframe=minesblue!40, boxrule=0.5pt, arc=2pt, left=6pt, right=6pt, top=3pt, bottom=3pt, fonttitle=\small\bfseries\color{minesblue}, title={Before You Start: Module 4}]
\small
\begin{enumerate}[nosep,leftmargin=*]
\item Read MRD Chapter~7 (Analog vs.\ Digital Systems) and Chapter~11 (DSP Fundamentals).
\item Calculate the LSB and temperature resolution for the NI~USB-6001 with your sensor before lab.
\item Prepare an RC anti-alias filter: you will need a resistor and capacitor from the lab stock.
\end{enumerate}
\end{tcolorbox}

\begin{warnbox}
\textbf{Aliasing is the single most common data acquisition mistake.} If you sample at 1\,kS/s without an anti-alias filter, any noise above 500\,Hz folds into your measurement as a false lower-frequency signal. The NI~USB-6001 has \textbf{no built-in anti-alias filter}; you must add an external RC or active filter before the analog input. The Master Reference Document, Chapter~11 (Digital Signal Processing), has the alias frequency formula and the ThermoRig switching-regulator example.
\end{warnbox}

\begin{tipbox}
Use $f_s \geq 10 \times f_{\max}$ in practice (not just $2\times$). This gives you a good waveform shape, margin for the anti-alias filter roll-off, and room for the FFT to resolve your signal clearly. For a 100\,Hz signal, sample at 1\,kS/s minimum.
\end{tipbox}

\begin{faqbox}[Q: How many bits do I actually need?]
Calculate the LSB: $\text{LSB} = V_{\text{range}} / 2^n$. Compare this to your sensor's output change per unit measurand. If your thermocouple produces 40\,$\mu$V/\textdegree{}C and your ADC LSB is 2.44\,mV (12-bit, 10\,V range), your temperature resolution is:
\[
\text{Temperature resolution} = \frac{\text{LSB}}{\text{Sensitivity}} = \frac{2.44\,\text{mV}}{0.040\,\text{mV/\textdegree{}C}} = 61\,\text{\textdegree{}C}
\]
That is useless. You need amplification (Chapter~15 (Amplifiers) of the MRD) to bring the signal into the ADC's usable range. With a gain of 100, the effective LSB becomes 24.4\,$\mu$V, giving 0.61\,\textdegree{}C resolution.
\end{faqbox}

%=============================================================================
\section{Module 5: Frequency Analysis (Fourier and FFT)}
%=============================================================================

\begin{priorities}
\begin{enumerate}[leftmargin=*, itemsep=1pt, topsep=0pt]
    \item Understand that any signal is a sum of sinusoids (Fourier principle).
    \item Use the FFT in LabVIEW to identify frequency components.
    \item Know how to choose $N$, $f_s$, and window type.
\end{enumerate}
\end{priorities}

\begin{tcolorbox}[colback=minesblue!3, colframe=minesblue!40, boxrule=0.5pt, arc=2pt, left=6pt, right=6pt, top=3pt, bottom=3pt, fonttitle=\small\bfseries\color{minesblue}, title={Before You Start: Module 5}]
\small
\begin{enumerate}[nosep,leftmargin=*]
\item Read MRD Chapter~9 (Fourier Series) and Chapter~10 (FFT).
\item Open the LabVIEW FFT Power Spectrum VI example on Canvas and run it with a known signal.
\item Know your three decisions before lab: $f_s$, $N$, and window type.
\end{enumerate}
\end{tcolorbox}

\begin{faqbox}[Q: How do I choose FFT parameters in LabVIEW?]
Three decisions: (1) \textbf{Sampling rate} $f_s$: at least $2.56 \times$ the highest frequency of interest (the ``2.56 rule'': the factor 2.56 provides a 28\% guard band above Nyquist for anti-alias filter roll-off). (2) \textbf{Number of points} $N$: determines frequency resolution $\Delta f = f_s / N$. Use a power of 2 for computational efficiency. (3) \textbf{Window}: use Hanning for general analysis; flat-top for accurate amplitude of known tones; rectangular only if the signal is perfectly periodic in the window.
\end{faqbox}

\begin{warnbox}
\textbf{Spectral leakage makes your FFT look wrong.} If the signal frequency does not exactly land on an FFT bin, energy ``leaks'' into adjacent bins, creating a broad smear instead of a sharp peak. The fix is windowing (Hanning is the default). Without a window, a 100.5\,Hz tone sampled at 1\,kS/s with $N = 1000$ ($\Delta f = 1$\,Hz) looks like it has energy spread across 20+ bins. With a Hanning window, the peak sharpens dramatically.
\end{warnbox}

\begin{protip}
\textbf{Average multiple FFT blocks for clean spectra.} A single FFT of noisy data has high variance. In LabVIEW, use ``RMS Averaging'' in the FFT Power Spectrum VI. Start with 10 averages; use 100 for noisy signals. The noise floor drops by $10\log_{10}(N_{\text{avg}})$\,dB.
\end{protip}

%=============================================================================
\section{Module 6: Analog and Digital Filter Design}
%=============================================================================

\begin{priorities}
\begin{enumerate}[leftmargin=*, itemsep=1pt, topsep=0pt]
    \item Design a passive RC low-pass filter and calculate its cutoff frequency.
    \item Know the three response types: Butterworth (flat), Chebyshev (sharp), Bessel (linear phase).
    \item Implement a digital filter in LabVIEW.
\end{enumerate}
\end{priorities}

\begin{tcolorbox}[colback=minesblue!3, colframe=minesblue!40, boxrule=0.5pt, arc=2pt, left=6pt, right=6pt, top=3pt, bottom=3pt, fonttitle=\small\bfseries\color{minesblue}, title={Before You Start: Module 6}]
\small
\begin{enumerate}[nosep,leftmargin=*]
\item Read MRD Chapter~12 (Analog Filter Design) and Chapter~13 (Digital Filter Design).
\item Calculate an RC filter cutoff ($f_c = 1/2\pi RC$) for your anti-alias requirement.
\item Locate the Sallen-Key op-amp circuit in the MRD; you may build one in lab.
\end{enumerate}
\end{tcolorbox}

\begin{faqbox}[Q: Which filter type should I use?]
\textbf{Butterworth} is the default because it has a maximally flat magnitude response in the passband (no ripple). Use it unless you have a specific reason not to. \textbf{Chebyshev}: when you need a steeper cutoff with fewer components (accept passband ripple). \textbf{Bessel}: when preserving the shape of a pulse or transient is critical (linear phase). The Master Reference Document, Chapters~12--13, has design procedures for both analog and digital implementations.
\end{faqbox}

\begin{warnbox}
\textbf{A first-order RC filter is only $-$20\,dB/decade.} At 10$\times$ the cutoff frequency, it attenuates by only 20\,dB (10$\times$). For anti-aliasing, this is often insufficient. Use at least a 2nd-order filter (Sallen-Key active, $-$40\,dB/decade) or combine an analog anti-alias filter with digital filtering in LabVIEW.
\end{warnbox}

\begin{protip}
\textbf{Build the analog filter first, then refine digitally.} Use a simple RC or 2nd-order Sallen-Key as the mandatory anti-alias filter before the ADC, then do the precision filtering in LabVIEW software. This gives you the best of both worlds: aliasing protection from analog hardware, and flexible/precise filtering from digital software.
\end{protip}

%=============================================================================
\section{Module 7: Signal Conditioning and Amplifiers}
%=============================================================================

\begin{priorities}
\begin{enumerate}[leftmargin=*, itemsep=1pt, topsep=0pt]
    \item Know the non-inverting, differential, and instrumentation amplifier configurations.
    \item Calculate the gain needed to match your sensor output to the ADC input range.
    \item Understand common-mode rejection and why it matters.
\end{enumerate}
\end{priorities}

\begin{tcolorbox}[colback=minesblue!3, colframe=minesblue!40, boxrule=0.5pt, arc=2pt, left=6pt, right=6pt, top=3pt, bottom=3pt, fonttitle=\small\bfseries\color{minesblue}, title={Before You Start: Module 7}]
\small
\begin{enumerate}[nosep,leftmargin=*]
\item Read MRD Chapter~15 (Amplifiers). Pay close attention to Section~15.7 (``The Brutal Realities'').
\item Calculate the gain you need: $G = V_{\text{ADC range}} / V_{\text{sensor max}}$.
\item Locate the INA128 and TL072 op-amps in the lab component drawers.
\end{enumerate}
\end{tcolorbox}

\begin{faqbox}[Q: Why can't I just connect the sensor directly to the ADC?]
Three reasons: (1) The sensor output is too small (thermocouples produce $\mu$V; the ADC needs volts). (2) The sensor impedance is too high (the ADC loads it, changing the reading). (3) Noise rides on the signal (the amplifier rejects common-mode noise). The Master Reference Document, Chapter~15 (Amplifiers), covers the complete signal conditioning chain from bridge circuit through instrumentation amplifier to ADC.
\end{faqbox}

\begin{warnbox}
\textbf{Loading error from mismatched impedances.} If your sensor output impedance is 10\,k$\Omega$ and you connect it to a 10\,k$\Omega$ input, you lose 50\% of the signal. The 100:1 rule: the measuring instrument's input impedance must be $\geq$100$\times$ the source impedance. Use a buffer amplifier (voltage follower) to solve this.
\end{warnbox}

\begin{protip}
\textbf{The INA128 is your best friend for bridge circuits.} It has $>$10\,G$\Omega$ balanced input impedance, $>$100\,dB CMRR, and gain set by a single resistor with $G = 1 + 49.4\,\text{k}\Omega/R_G$. For $G=100$: $R_G = 49.4\,\text{k}\Omega/99 = 499\,\Omega$ (use 499\,$\Omega$ standard). Whenever you measure a Wheatstone bridge (strain gauge, RTD, load cell), reach for an instrumentation amplifier first.
\end{protip}

%=============================================================================
\section{Module 8: High-Power Devices and Switching}
%=============================================================================

\begin{priorities}
\begin{enumerate}[leftmargin=*, itemsep=1pt, topsep=0pt]
    \item Know when to use a MOSFET, relay, or SSR.
    \item Always add a flyback diode for inductive loads.
    \item Understand PWM for proportional power control.
\end{enumerate}
\end{priorities}

\begin{tcolorbox}[colback=minesblue!3, colframe=minesblue!40, boxrule=0.5pt, arc=2pt, left=6pt, right=6pt, top=3pt, bottom=3pt, fonttitle=\small\bfseries\color{minesblue}, title={Before You Start: Module 8}]
\small
\begin{enumerate}[nosep,leftmargin=*]
\item Read MRD Chapter~16 (High-Power Devices and Switches).
\item Identify your actuator type and locate the appropriate driver (MOSFET, SSR, or H-bridge) in the lab stock.
\item Always have a 1N5819 Schottky diode ready for any inductive load.
\end{enumerate}
\end{tcolorbox}

\begin{warnbox}
\textbf{Inductive loads (motors, solenoids, relays) will destroy your MOSFET without a flyback diode.} When you switch off an inductive load, the collapsing magnetic field generates a voltage spike of hundreds or thousands of volts. A 1N5819 Schottky diode across the load (cathode to positive) costs \$0.10 and prevents \$5--\$50 of damage. See the Master Reference Document, Chapter~16 (High-Power Devices), for the spike calculation and circuit diagram.
\end{warnbox}

\begin{tipbox}
\textbf{For DC loads}: use an N-channel logic-level MOSFET (IRLZ44N) in the low side. For AC loads: use a solid-state relay (SSR). For bidirectional motor control: use an H-bridge IC (L298N, DRV8871). Never switch AC mains with a bare MOSFET.
\end{tipbox}

%=============================================================================
\section{Module 8B: Sensors and Actuators Fundamentals}
%=============================================================================

\begin{priorities}
\begin{enumerate}[leftmargin=*, itemsep=1pt, topsep=0pt]
    \item Understand transduction: how sensors convert physical quantities to electrical signals.
    \item Know the key specs: range, sensitivity, accuracy, bandwidth, output impedance.
    \item Match actuator type to load requirements before selecting a driver circuit.
\end{enumerate}
\end{priorities}

\begin{faqbox}[Q: How do I choose between a thermocouple, RTD, and thermistor?]
\textbf{Thermocouple}: widest range ($-200$ to $1250$\,\textdegree{}C for Type K), fastest response, self-generating (no excitation needed), but lowest output ($\mu$V-level, needs amplifier). \textbf{RTD (Pt100)}: best accuracy ($\pm$0.1\,\textdegree{}C), excellent linearity, but slower and requires bridge excitation. \textbf{Thermistor}: highest sensitivity (10$\times$ TC), but highly nonlinear and limited range ($-40$ to $125$\,\textdegree{}C). Pick based on your range and accuracy requirements first, then consider signal conditioning complexity. Chapter~6 (Sensors and Actuators) of the MRD has the full selection table.
\end{faqbox}

\begin{warnbox}
\textbf{Every actuator needs a driver circuit.} No motor, solenoid, or heater connects directly to a microcontroller GPIO. Motors need H-bridges or ESCs. Solenoids need MOSFETs with flyback diodes. Heaters need SSRs or MOSFETs with current limiting. Forgetting the driver is the most common actuator integration mistake.
\end{warnbox}

\begin{protip}
\textbf{Check sensor output compatibility before ordering.} Verify three things: (1) the output voltage range fits your ADC input range (a 10\,V sensor on a 3.3\,V ADC will saturate and may damage the ADC), (2) the output impedance is compatible ($>$100:1 ratio with your ADC input impedance), and (3) the response time is fast enough for your signal ($\tau_{\text{sensor}} < 1/(5 f_{\text{signal}})$). Chapter~6 (Sensors and Actuators) has the sensor selection decision framework.
\end{protip}

%=============================================================================
\section{Module 8C: Fluid Power Fundamentals}
%=============================================================================

\begin{priorities}
\begin{enumerate}[leftmargin=*, itemsep=1pt, topsep=0pt]
    \item Understand pressure, flow, and their relationship (Bernoulli, continuity).
    \item Know how differential-pressure flow measurement works (orifice, Venturi, Pitot).
    \item Distinguish hydraulic vs. pneumatic system tradeoffs.
\end{enumerate}
\end{priorities}

\begin{warnbox}
\textbf{This module is project-dependent.} Only read this section if your project involves pumps, valves, flow measurement, or hydraulic/pneumatic actuation.
\end{warnbox}

\begin{faqbox}[Q: My project has a pump or flow control. Where do I start?]
Start with Chapter~22 (Fluid Power Fundamentals) of the MRD. Determine your flow rate and pressure requirements. For precise metering of small flows ($<$5\,L/min), use a peristaltic or gear pump. For high-flow cooling or mixing ($>$10\,L/min), use a centrifugal pump. For pneumatic actuation, you need an FRL (filter-regulator-lubricator) unit upstream of any pneumatic cylinder or valve.
\end{faqbox}

\begin{warnbox}
\textbf{Never check for hydraulic leaks with your hand.} A pinhole leak at 200\,bar produces a jet that can penetrate skin and inject fluid into tissue (fluid injection injury, a surgical emergency). Use a piece of cardboard to locate leaks. Always relieve pressure before disconnecting fittings.
\end{warnbox}

\begin{tipbox}
\textbf{Flow measurement shortcut}: for clean water or air in pipes under 50\,mm, a small turbine flowmeter with pulse output is the easiest to interface with the ESP32. Count pulses with a hardware timer and convert to flow rate using the meter's K-factor (pulses per liter). The ThermoRig cooling loop example in Chapter~22 walks through the complete calculation.
\end{tipbox}

%=============================================================================
\section{Module 8D: Systematic Debugging and Troubleshooting}
%=============================================================================

\begin{priorities}
\begin{enumerate}[leftmargin=*, itemsep=1pt, topsep=0pt]
    \item Always debug bottom-up: component $\to$ interface $\to$ subsystem $\to$ system.
    \item Use the five-step protocol: reproduce, isolate, measure, hypothesize, fix and regression test.
    \item Change only one variable between tests.
\end{enumerate}
\end{priorities}

\begin{faqbox}[Q: I wired everything together and nothing works. Where do I start?]
Disconnect everything. Go back to Level~1 of the debugging pyramid. Verify each component in isolation: does the sensor output the right voltage? Does the amplifier have the correct gain with a signal generator input? Does the MCU blink an LED? Does the motor spin with a direct battery connection? Only after every component passes individually do you begin reconnecting them, one interface at a time, testing after each connection. This feels slower, but it is 5--10$\times$ faster than staring at a fully assembled system that does not work. See Chapter~17 (Systematic Debugging) of the MRD for the full methodology.
\end{faqbox}

\begin{warnbox}
\textbf{``I changed three things and now it works'' is not debugging.} If you swapped a component, changed the wiring, and edited the firmware simultaneously, you do not know which change fixed the problem. Worse, one of the other changes may have introduced a new latent fault. Always change one variable at a time, test, and record the result before making the next change.
\end{warnbox}

\begin{tipbox}
\textbf{The Signal Walk is your most powerful debugging technique.} Start at the sensor with an oscilloscope probe. Measure the output. Move one stage downstream (after the divider, after the amplifier, at the ADC input). At each point, compare measured vs. expected. The first point where they diverge localizes the fault to that stage. For example: measure thermocouple output ($\sim$1\,mV), then INA output ($\sim$100\,mV at $G=100$), then filter output ($\sim$100\,mV, less noise), then ADC input pin. The first point where measured $\neq$ expected localizes the fault. Five minutes of signal walking saves hours of guessing.
\end{tipbox}

\begin{protip}
\textbf{Keep a debug log.} Date, symptom, tests performed, measurements taken, hypothesis, fix applied, result. When the same failure appears two weeks later (it will), you will have the answer in 30 seconds instead of 2 hours. The debug log also demonstrates professional engineering practice in your final report.
\end{protip}

\begin{faqbox}[Q: My sensor readings are fine until the motor turns on. Then everything gets noisy.]
Classic ground loop or EMI problem. The motor's switching current is coupling into the sensor signal path. Three fixes, in order of effectiveness: (1)~Star grounding: separate ground wires for sensor and motor circuits, joined at one point near the supply. (2)~Add a low-pass filter at the amplifier input to reject high-frequency motor noise. (3)~Route sensor wires physically away from motor power wires and use twisted-pair for the sensor leads. Chapter~17 (Systematic Debugging) covers the full integration debugging protocol.
\end{faqbox}

%=============================================================================
\section{Module 9: Control Systems and PID}
%=============================================================================

\begin{priorities}
\begin{enumerate}[leftmargin=*, itemsep=1pt, topsep=0pt]
    \item Understand the difference between open-loop and closed-loop control.
    \item Know what each PID term does: P (speed), I (accuracy), D (damping).
    \item Tune a PID controller using the Ziegler-Nichols or manual method.
\end{enumerate}
\end{priorities}

\begin{tcolorbox}[colback=minesblue!3, colframe=minesblue!40, boxrule=0.5pt, arc=2pt, left=6pt, right=6pt, top=3pt, bottom=3pt, fonttitle=\small\bfseries\color{minesblue}, title={Before You Start: Module 9}]
\small
\begin{enumerate}[nosep,leftmargin=*]
\item Read MRD Chapter~18 (Open-Loop), Chapter~19 (Closed-Loop), and Chapter~20 (PID Control).
\item Open the LabVIEW PID VI example on Canvas and familiarize yourself with the gain inputs.
\item Review the manual tuning sequence: P first, then I, then D.
\end{enumerate}
\end{tcolorbox}

\begin{faqbox}[Q: My P-only controller has a steady-state offset. Why?]
P-only control requires $e > 0$ to produce $u > 0$. If $e = 0$, the controller output is zero, and the process drifts. The system settles where the disturbance (heat loss, friction) exactly balances the control effort at some nonzero error. Only integral (I) action can drive the error to zero, because the integral accumulates error over time and keeps increasing the output until $e = 0$. See the Master Reference Document, Chapter~20 (PID Control), for the detailed explanation.
\end{faqbox}

\begin{warnbox}
\textbf{Integral windup is the \#1 PID debugging headache.} If your system overshoots massively after a large setpoint change, the integral term has ``wound up'' while the actuator was saturated. Enable anti-windup: stop accumulating the integral when the output hits its limit. The Master Reference Document, Chapter~20 (PID Control), has the discrete-time pseudocode with anti-windup clamping.
\end{warnbox}

\begin{protip}
\textbf{Manual PID tuning sequence}: (1)~Start P-only; increase $K_p$ until the response is fast with $\sim$20\% overshoot. (2)~Add I; start $K_i$ small, increase until the steady-state error vanishes. If overshoot grows, reduce $K_i$. (3)~Add D; start $K_d$ small, increase to reduce overshoot. If the output becomes noisy, $K_d$ is too high. Tune $K_p$ first, then $K_i$, then $K_d$; never all three at once. If the output becomes noisy when you add D, reduce $K_d$ or add a derivative filter (the LabVIEW PID VI has parameter $N$; set to 10--20).
\end{protip}

\begin{faqbox}[Q: My PID output is noisy and jerky. What is wrong?]
The derivative term amplifies noise. Three fixes: (1) reduce $K_d$, (2) add a low-pass filter on the derivative (the LabVIEW PID VI has a built-in derivative filter parameter $N$; set it to 10--20), (3) filter the process variable before it enters the PID (a moving average of 5--10 points works well for thermocouples and other slow sensors).
\end{faqbox}

%=============================================================================
\section{LabVIEW Survival Guide}
%=============================================================================

\begin{faqbox}[Q: My VI runs but the DAQ gives ``resource reserved'' errors. What happened?]
You forgot to wire \texttt{DAQmx Clear Task} outside your While Loop. When LabVIEW stops (even by pressing Stop), the DAQ task stays locked unless explicitly cleared. Always wire Clear Task on the exit path of the loop, connected via the error wire chain.
\end{faqbox}

\begin{warnbox}
\textbf{Six LabVIEW mistakes that cost 80\% of grading deductions} (definitive list; see MRD LabVIEW appendix for detailed explanations):
\begin{enumerate}[nosep,leftmargin=*]
\item No loop timing (Wait function missing; loop runs at max speed).
\item DAQ task not cleared (``resource reserved'' on next run).
\item Wrong channel mode (RSE instead of Differential).
\item Chart vs.\ Graph confusion (Chart for real-time; Graph for post-acquisition).
\item PID loop rate not controlled (gains become unpredictable).
\item No anti-alias filter before ADC.
\end{enumerate}
Eliminate all six and you avoid most lost points.
\end{warnbox}

\begin{tipbox}
\textbf{Debug with Highlight Execution and Probes.} The light bulb icon animates dataflow so you can see values moving along wires. Right-click any wire and add a Probe to watch its value in real time. These two tools solve 90\% of LabVIEW bugs.
\end{tipbox}

\begin{protip}
\textbf{Log data inside the loop, not after.} Write each data point to a file within the While Loop using ``Write to Measurement File.'' If LabVIEW crashes (or you press Stop too early), you still have all the data collected up to that point. Accumulating data in arrays and writing at the end risks losing everything.
\end{protip}

%=============================================================================
\section{Design Challenge Tips}
%=============================================================================

\begin{faqbox}[Q: How should I approach a Design Challenge?]
Design Challenges are timed, in-class engineering problems. The key is structured thinking under pressure:
\begin{enumerate}[leftmargin=*, itemsep=2pt, topsep=0pt]
    \item \textbf{Read the entire prompt first} (2 minutes). Identify what you are measuring, what accuracy is required, and what hardware is available.
    \item \textbf{Sketch the signal chain} (3 minutes). Draw the five blocks: sensor $\to$ conditioning $\to$ DAQ $\to$ processing $\to$ output. Identify what each block does.
    \item \textbf{Calculate before building} (5 minutes). What gain do you need? What filter cutoff? What sampling rate? What is the expected output voltage range?
    \item \textbf{Build and test incrementally} (remaining time). Wire the sensor first and verify the raw signal. Then add amplification. Then add the filter. Then set up LabVIEW. Test after each step.
\end{enumerate}
\end{faqbox}

\begin{warnbox}
\textbf{Do not start wiring until you have done the math.} The most common Design Challenge failure is connecting everything, running LabVIEW, and seeing garbage on the screen with no idea why. Five minutes of calculation (gain, filter, sampling rate) prevents two hours of blind debugging.
\end{warnbox}

\begin{protip}
\textbf{Keep a cheat sheet of key formulas.} You will want these at your fingertips during Design Challenges:
\begin{itemize}[nosep,leftmargin=*]
\item $f_c = 1/(2\pi RC)$ (RC filter cutoff)
\item $\text{LSB} = V_{\text{range}}/2^n$ (ADC resolution)
\item $G = 1 + R_f/R_g$ (non-inverting amp gain)
\item $f_s > 2 f_{\max}$ (Nyquist), practical $f_s \geq 10 f_{\max}$
\item $\text{SNR}_{\text{dB}} = 20\log_{10}(V_s/V_n)$
\item $e_{ss} = 0$ requires integral action (PI or PID)
\end{itemize}
The Quick Reference Tables in the Master Reference Document consolidate all of these.
\end{protip}

\begin{faqbox}[Quick Reference: Common Component Values]
\small
\begin{tabularx}{\textwidth}{@{} l X @{}}
\toprule
\textbf{Component} & \textbf{Key Values} \\
\midrule
Type K thermocouple & Sensitivity: $\sim$41\,$\mu$V/\textdegree{}C; Range: $-200$ to $1250$\,\textdegree{}C \\
Pt100 RTD & 100\,$\Omega$ at 0\,\textdegree{}C; TCR: 0.385\,$\Omega$/\textdegree{}C \\
NI USB-6001 & $\pm$10\,V input; 14-bit; 20\,kS/s; 8 AI (4 diff); no built-in anti-alias filter \\
INA128 & $G = 1 + 49.4\,\text{k}\Omega/R_G$; CMRR $>$ 100\,dB; Input $Z > 10$\,G$\Omega$ \\
IRLZ44N MOSFET & Logic-level gate; $R_{DS(on)} = 0.022\,\Omega$; $I_D = 47$\,A; $V_{DS} = 55$\,V \\
1N5819 Schottky & $V_f = 0.3$\,V; $I_{\max} = 1$\,A; fast recovery for flyback protection \\
\bottomrule
\end{tabularx}
\end{faqbox}

%=============================================================================
\section{Lab Practical Exam Preparation}
%=============================================================================

\begin{faqbox}[Q: What should I study for the lab practical exams?]
Lab practicals test both theory and hands-on skills. You may be asked to: wire a circuit on a breadboard, configure the NI~DAQ, write a LabVIEW VI from scratch, perform a calibration, use the FFT, design a filter, or tune a PID controller. The best preparation is to re-do each lab exercise from scratch without looking at your previous work. If you can do it from memory, you are ready.
\end{faqbox}

\begin{tipbox}
\textbf{Practice the physical skills.} Reading a resistor color code, wiring a Wheatstone bridge on a breadboard, connecting the NI~DAQ terminals, and using a multimeter and bench power supply are all skills that improve dramatically with practice. Do not wait until the exam to discover you cannot remember which DAQ terminal is AI0+.
\end{tipbox}

\begin{warnbox}
\textbf{Common lab practical mistakes (avoid these):}
\begin{enumerate}[nosep,leftmargin=*]
\item Wiring the DMM current jacks across a voltage source (blows the fuse).
\item Forgetting the DAQmx Clear Task (``resource reserved'' on the next run).
\item Using RSE mode instead of Differential (doubles the noise).
\item No anti-alias filter before the ADC (aliased data looks correct until you check the FFT).
\item Connecting the oscilloscope ground clip to a non-ground node (creates a short through earth ground).
\item Setting PID gains without first running P-only to understand the process response.
\end{enumerate}
\end{warnbox}

\begin{protip}
\textbf{Specific skills to practice:} (1)~Wire a non-inverting amplifier with a specified gain in under 5 minutes. (2)~Configure a DAQmx Read VI with correct channel, terminal config, and sample rate. (3)~Set up a bench supply with voltage and current limits. (4)~Measure V, I, and R with a DMM (current in series, voltage in parallel). (5)~Identify RC filter cutoff by sweeping a signal generator and measuring the $-3$\,dB point.
\end{protip}

%=============================================================================
\section{Quick Reference: Module Checklist}
%=============================================================================

\begin{table}[htbp]
\centering\small
\begin{tabularx}{\textwidth}{@{} c l X l @{}}
\toprule
\textbf{Wk} & \textbf{Module} & \textbf{Key Skills} & \textbf{MRD Ch.} \\
\midrule
1--2 & Measurement Fundamentals & Sensor terminology, calibration, bridge circuits & 1--2 \\
3--4 & Dynamic Meas.\ \& Error & Time constants, frequency response, uncertainty propagation & 3--4 \\
5 & Statistics \& Sensors & Experimental data analysis, sensor/actuator selection & 5--6 \\
6--7 & Signals \& Digitization & ADC specs, aliasing, SNR, anti-alias filters & 7--8, 11 \\
8--9 & Frequency Analysis & Fourier series, FFT, windowing, PSD & 9--10 \\
10--11 & Filtering & Analog (RC, Sallen-Key), digital (FIR/IIR), LabVIEW filters & 12--13 \\
12 & Signal Conditioning & Amplifiers, INA, circuits, high-power switches, PWM & 14--16 \\
12--13 & Debugging & Bottom-up methodology, signal walk, integration protocol & 17 \\
13--14 & Control Systems & Open/closed-loop, PID tuning, LabVIEW PID VI & 18--20 \\
15 & Engineering Science & Heat transfer, fluid power (if applicable) & 21--22 \\
16 & Integration \& Exams & Complete system design, lab practicals & All \\
\bottomrule
\end{tabularx}
\end{table}

\vspace{1cm}
\begin{center}
\textcolor{minesblue}{\rule{0.4\textwidth}{1pt}}\\[6pt]
\textit{Measure carefully, filter thoughtfully, control precisely, and document everything.}\\[6pt]
\textcolor{minesblue}{\rule{0.4\textwidth}{1pt}}
\end{center}

\end{document}
